% !TeX root = ../../Book3.tex
\section{零点序列}
\label{b3:sec:1903}

Jensen 公式限制零点的数量，但还可以把整个零点序列压成一个收敛指数。下一章构造典范乘积时，这个指数决定必须扣掉多少低阶项。

\subsection{收敛指数}
\label{b3:subsec:190301}

按重数列出趋于无穷的非零零点 \(a_n\)。定义\term[lingdian-shoulian-zhishu]{零点收敛指数}[exponent of convergence]为 \(\lambda=\inf\{s>0:\sum|a_n|^{-s}<\infty\}\)。有限零点序列取零。对于无穷序列，若 \(n(r)\) 为计数，则
\[
\lambda=\limsup_{r\to\infty}\frac{\log n(r)}{\log r}.
\]
证明可用分部求和：\(\sum_{|a_n|>1}|a_n|^{-s}=s\int_1^\infty n(t)t^{-s-1}\,\mathrm{d}t\) 加有限边界项。若计数为 \(O(t^{s-\varepsilon})\) 积分收敛；若序列幂和收敛，则 \(n(r)\le r^s\sum|a_n|^{-s}\)，给反向。

\subsection{零点密度}
\label{b3:subsec:190302}

\begin{proposition}\label{b3:prop:19-zero-density}
有限阶 \(\rho\) 的非恒零整函数具有 \(n(r)=O_\varepsilon(r^{\rho+\varepsilon})\)，零点收敛指数不超过 \(\rho\)。
\end{proposition}
\begin{proof}
先除去原点零点。Jensen 估计取外半径 \(2r\)，得到 \(n(r)\log2\le\log M(2r)-\log|f(0)|\)。阶的定义给计数界；取稍小的 \(\varepsilon\) 后，上一积分判据使每个 \(s>\rho\) 的零点幂和收敛。
\end{proof}
原点的有限重数只增加常数项，不改变收敛指数。

\subsection{增长限制}
\label{b3:subsec:190303}

反向并不唯一：\(e^{z^m}\) 无零点而阶 \(m\)，所以零点只能约束增长的一部分。\(\sin\pi z\) 的零点为整数，收敛指数一，和阶一相同。

若指定 \(a_n=n^\alpha\)，其收敛指数为 \(1/\alpha\)。任何具有这些零点的整函数必须有阶至少 \(1/\alpha\)。下一章将证明离散零点总能实现，但若还要求有限阶，就要检查这个收敛指数。


\ProseParagraphBreak

零点计数的幂界可用二进尺度再证明收敛指数判据。若 \(n(r)\le Cr^a\)，在 \(2^j<|a_n|\le2^{j+1}\) 内，\(s>a\) 时幂和至多 \(C2^{(j+1)a-js}\)，对 \(j\) 求和为收敛几何级数。若 \(\sum|a_n|^{-s}<\infty\)，则半径 \(r\) 内每项至少 \(r^{-s}\)，故 \(n(r)r^{-s}\) 有界。这个分组证明不要求计数函数可微，也自然保留重数；分部求和则给更精确的常数与临界积分。

\ProseParagraphBreak

在临界指数上，下确界不决定幂和是否收敛。零点模长为 \(n\) 时收敛指数一，临界一次幂和发散；模长为 \(n\log^2(n+1)\) 时指数仍一，却临界一次幂和收敛。若模长为 \(n\log(n+1)\)，临界再次发散。对数修正不会改变 \(\log n(r)/\log r\) 的上极限，却决定下一章的最小初等因子次数。阶只决定粗尺度，属还要检查临界幂和，二者不可互换。

\ProseParagraphBreak

零点重数也能提高收敛指数。若每个正整数 \(n\) 的重数为 \(n\)，幂和变成 \(\sum n^{1-s}\)，收敛指数为二；计数约为 \(r^2/2\)，因此任何实现函数阶至少二。若仅按不同零点个数，则会得到错误下界一。对整函数增长、典范乘积和 Jensen 的贡献，统一按重数计数是必要条件，并非一个可选的记账方式。

\ProseParagraphBreak

增长下界只来自零点径向分布，角位置可以有进一步作用。例如整数零点在两条相反射线上排列，正弦给出阶一实现；把正整数零点全部放在同一射线，最小属仍是一，但沿某些方向的精确增长会不同。下一章的乘积修正会显式出现一次项，说明方向排列影响低次指数自由度。当前的 \(\lambda\le\rho\) 是普适必要条件，还没有声称它给所有方向的精确渐近。

\ProseParagraphBreak

若不仅知道阶，还知道正常型，Jensen 估计可给更具体的密度常数。假设 \(0<\rho<\infty\)、\(\tau<\infty\)，先去掉原点零点。对任意固定 \(u>1\)，有
\[
 n(r)\log u\le\log M(ur)-\log|f(0)|.
\]
除以 \(r^\rho\) 后取上极限，右侧不超过 \(\tau u^\rho\)。函数 \(u^\rho/\log u\) 在 \(\log u=1/\rho\) 处取最小值 \(e\rho\)，因此
\[
 \limsup_{r\to\infty}\frac{n(r)}{r^\rho}\le e\rho\tau.
\]
这仍是一般上界，不声称常数对每类函数都最优。它至少说明最小型 \(\tau=0\) 时 \(n(r)=o(r^\rho)\)，比只给 \(O(r^{\rho+\varepsilon})\) 更细。有限个原点零点除以 \(r^\rho\) 后消失，不改变结果。
\ProseParagraphBreak
对正弦函数，\(n(r)\) 与 \(r\) 同阶；对指数函数，计数恒为零，二者却有相同的阶与型。这再次说明增长和零点不是互相决定的两组资料。若只给定零点序列，下一章构造的典范乘积会固定其中一部分增长，而允许的无零点因子仍可改变最大模。使用密度估计时，正确方向是“增长上界限制零点”，不能不加额外假设把它倒置为“零点上界限制函数”。
